\section{Introduction}\label{sec:introduction}

Suppose that a finite carrier algebra $A$ and a finite operator algebra $B$
are known completely, but only some values of an operation
$\alpha:B\times A\to A$ have been prescribed. Requiring each named operator
to preserve the native operations of $A$ connects the missing cells to one
another. It can also force originally distinct carrier elements to become
equal. A semantics that immediately fills every blank with a carrier element
can therefore obscure what the original equations actually determine.

This paper separates three questions. Which equalities among the original
carrier names hold in every model? Which missing action values already equal
carrier names in every model? On which carrier quotients can all missing
values be chosen inside the quotient? These are, respectively, questions
about \emph{protection}, \emph{determinacy}, and \emph{completion}.
They have different answers, even for a single row on a four-element algebra.

We use initial-algebra semantics for a finite ground presentation. The native
tables specify named copies of $A$ and $B$; compatibility is required only at
named inputs. An unspecified cell denotes a total operation value in every
model, but that value need not belong to the named carrier image. We call
such an initial value \emph{external}. This use of external values is
essential to distinguish necessary identifications from choices introduced
by a completion.

The main results are as follows.
\begin{enumerate}
\item \textbf{Exact forced quotient.} Each row is a homomorphic image of the
carrier, glued to it along the supplied data. A row pushout preserves the
carrier exactly when the data induce a homomorphism $h:S\to Q$ whose kernel
extends to $Q$ without acquiring new pairs inside $S$
(Theorem~\ref{thm:raw}). The least quotient with this property for every row
is the exact carrier quotient of the initial algebra
(Theorem~\ref{thm:leaststable}). A monotone feedback operator computes it in
at most $|A|-1$ strict quotient steps, and this bound is sharp.
\item \textbf{Exact forced domain.} For a protected row, the ambient
congruence $\kappa=\Cg_Q(\ker h)$ is exactly the row kernel. A point has a
forced carrier value precisely when its $\kappa$-class meets $S$
(Theorem~\ref{thm:kernelsat}). Theorem~\ref{thm:globalrows} transfers this
description to all rows in the full initial algebra. The resulting
description also gives exact class counts on the named interface at proof
horizons one and two (Theorem~\ref{thm:interface}).
\item \textbf{Completion and repair.} Completing a row on $Q$ is equivalent
to splitting its carrier embedding into the pushout
(Theorem~\ref{thm:splitting}). The repair congruences are exactly the carrier
kernels of models whose named action values stay inside the named carrier
image (Theorem~\ref{thm:nojunk}). Examples separate this model restriction
from unrestricted equational consequence.
\end{enumerate}

The categorical and logical tools are classical: initial algebras,
amalgamation, congruence generation, and ground congruence closure. The
contribution is their exact combination for arbitrary incomplete action
rows, including inconsistent supplied values. In particular, the input is
not assumed to be a partial action or even a partial function. The
homomorphic partial maps arise after the required quotient has been found.

This article is the semantic foundation of a three-paper programme. The
companion manuscript on equality-visibility depth~\cite{ShcherbakovII}
studies the dependence of finite-ground equality on a chosen presentation.
The planned third paper concerns the computational core and checkable
certificates. Here we retain only the finite-ground filtration needed to
connect the structural results to those later questions. All results used
in this paper are proved here or attributed to the preceding literature;
the argument does not depend on either companion being published.

Sections~\ref{sec:setup}--\ref{sec:mainquotient} develop the pure-row theory.
Section~\ref{sec:examples} gives the examples and named-interface filtration.
Section~\ref{sec:repair} treats completion and repair.
Sections~\ref{sec:actionlaws} and~\ref{sec:onesorted} delimit the effect of
additional action laws and one-sorted coupling.
Section~\ref{sec:computation} explains finite verification, and
Section~\ref{sec:related} situates the results in the literature.
